% Options for packages loaded elsewhere
\PassOptionsToPackage{unicode}{hyperref}
\PassOptionsToPackage{hyphens}{url}
\documentclass[
]{article}
\usepackage{xcolor}
\usepackage{amsmath,amssymb}
\setcounter{secnumdepth}{-\maxdimen} % remove section numbering
\usepackage{iftex}
\ifPDFTeX
  \usepackage[T1]{fontenc}
  \usepackage[utf8]{inputenc}
  \usepackage{textcomp} % provide euro and other symbols
\else % if luatex or xetex
  \usepackage{unicode-math} % this also loads fontspec
  \defaultfontfeatures{Scale=MatchLowercase}
  \defaultfontfeatures[\rmfamily]{Ligatures=TeX,Scale=1}
\fi
\usepackage{lmodern}
\ifPDFTeX\else
  % xetex/luatex font selection
\fi
% Use upquote if available, for straight quotes in verbatim environments
\IfFileExists{upquote.sty}{\usepackage{upquote}}{}
\IfFileExists{microtype.sty}{% use microtype if available
  \usepackage[]{microtype}
  \UseMicrotypeSet[protrusion]{basicmath} % disable protrusion for tt fonts
}{}
\makeatletter
\@ifundefined{KOMAClassName}{% if non-KOMA class
  \IfFileExists{parskip.sty}{%
    \usepackage{parskip}
  }{% else
    \setlength{\parindent}{0pt}
    \setlength{\parskip}{6pt plus 2pt minus 1pt}}
}{% if KOMA class
  \KOMAoptions{parskip=half}}
\makeatother
\usepackage{longtable,booktabs,array}
\usepackage{calc} % for calculating minipage widths
% Correct order of tables after \paragraph or \subparagraph
\usepackage{etoolbox}
\makeatletter
\patchcmd\longtable{\par}{\if@noskipsec\mbox{}\fi\par}{}{}
\makeatother
% Allow footnotes in longtable head/foot
\IfFileExists{footnotehyper.sty}{\usepackage{footnotehyper}}{\usepackage{footnote}}
\makesavenoteenv{longtable}
\setlength{\emergencystretch}{3em} % prevent overfull lines
\providecommand{\tightlist}{%
  \setlength{\itemsep}{0pt}\setlength{\parskip}{0pt}}
\usepackage{bookmark}
\IfFileExists{xurl.sty}{\usepackage{xurl}}{} % add URL line breaks if available
\urlstyle{same}
\hypersetup{
  pdftitle={Problem set 1},
  hidelinks,
  pdfcreator={LaTeX via pandoc}}

\title{Problem set 1}
\author{}
\date{2026-01-16}

\begin{document}
\maketitle

\subsection{Welcome}\label{welcome}

Welcome to Problem Set 1! There are some ``problem'' exercises and one
extended Stata exercise. You'll need to submit two things on
Brightspace: a problem set and log file. If you have trouble with the
Stata basics, head back to \href{../01-lab/}{Lab 1}.

\begin{quote}
\emph{Tip: If, after doing these problems, you still want more practice,
the odd-numbered Exercises (not Empirical Exercises) in Chapters 2 and 3
of Stock and Watson are quite useful, and solutions are available
online.}
\end{quote}

Note that there are superscripted numbers\footnote{See?! Neat! :)}
throughout the page that provide additional information/suggestions to
help you out.

\subsection{Exercises}\label{exercises}

{\def\LTcaptype{} % do not increment counter
\begin{longtable}[]{@{}rrrrr@{}}
\toprule\noalign{}
\endhead
\bottomrule\noalign{}
\endlastfoot
X & -1 & 0 & 1 & 4 \\
\(P(X=x)\) & 0.25 & 0.30 & 0.40 & 0.05 \\
\end{longtable}
}

\begin{enumerate}
\def\labelenumi{\arabic{enumi}.}
\item
  Consider the above random variable, \(X\), with its associated
  probability distribution:

  \begin{enumerate}
  \def\labelenumii{\alph{enumii}.}
  \item
    \emph{Draw} the probability distribution function and the cumulative
    distribution function. (That is, you should make a figure/graph!)
  \item
    What is the expected value of X? That is, what is \(E[X]\)?
  \item
    What is the variance of X?
  \end{enumerate}
\item
  The following table gives the joint probability distribution between
  employment status and college graduation among those either employed
  or looking for work (unemployed) in the working-age U.S. population
  for 2017:

  {\def\LTcaptype{} % do not increment counter
  \begin{longtable}[]{@{}
    >{\raggedright\arraybackslash}p{(\linewidth - 6\tabcolsep) * \real{0.4154}}
    >{\raggedright\arraybackslash}p{(\linewidth - 6\tabcolsep) * \real{0.2769}}
    >{\raggedright\arraybackslash}p{(\linewidth - 6\tabcolsep) * \real{0.2308}}
    >{\raggedright\arraybackslash}p{(\linewidth - 6\tabcolsep) * \real{0.0769}}@{}}
  \toprule\noalign{}
  \begin{minipage}[b]{\linewidth}\raggedright
  \end{minipage} & \begin{minipage}[b]{\linewidth}\raggedright
  Unemployed (\(Y=0\))
  \end{minipage} & \begin{minipage}[b]{\linewidth}\raggedright
  Employed (\(Y=1\))
  \end{minipage} & \begin{minipage}[b]{\linewidth}\raggedright
  Total
  \end{minipage} \\
  \midrule\noalign{}
  \endhead
  \bottomrule\noalign{}
  \endlastfoot
  Non-college grads (\(X = 0\)) & 0.026 & 0.576 & 0.602 \\
  College grads (\(X = 1\)) & 0.009 & 0.389 & 0.398 \\
  \textbf{Total} & 0.035 & 0.965 & 1.000 \\
  \end{longtable}
  }

  \begin{enumerate}
  \def\labelenumii{\alph{enumii}.}
  \item
    Compute \(E[Y]\).
  \item
    The unemployment rate is the fraction of the labor force that is
    unemployed. Show that the unemployment rate is given by
    \(1 - E[Y]\).
  \item
    Calculate \(E[Y|X=1]\) and \(E[Y|X=0]\).
  \item
    Calculate the unemployment rate for college graduates and for
    non-college graduates
  \item
    A randomly selected member of this population reports being
    unemployed. What is the probability that this worker is a college
    graduate? A non-college graduate?
  \end{enumerate}

  \hspace{0pt} f.~Are educational achievement and employment status
  independent? Explain.
\item
  Compute the following probabilities\footnote{Remember that we
    conventionally write \(N(\mu,\sigma^2)\), so the second term is the
    \emph{variance}, not the standard deviation.}:

  \begin{enumerate}
  \def\labelenumii{\alph{enumii}.}
  \item
    If \(Y\) is distributed \(N(1,4)\), find \(Pr(Y \leq 3)\)
  \item
    If \(Y\) is distributed \(N(3,9)\), find \(Pr(Y >0 )\)
  \item
    If \(Y\) is distributed \(N(50,25)\), find \(Pr(40 \leq Y \leq 52)\)
  \item
    If \(Y\) is distributed \(N(5,2)\), find \(Pr(6 \leq Y \leq 8)\)
  \end{enumerate}
\item
  For a randomly selected county in the United States, let \(X\)
  represent the proportion of adults over age 65 who are employed (the
  elderly employment rate). Then, \(X\) is restricted to a value between
  zero and one. Suppose that the cumulative distribution function for
  \(X\) is given by \(F(x) = 3x^2 - 2x^3\) for \(0 \leq x \leq 1\).

  \begin{enumerate}
  \def\labelenumii{\alph{enumii}.}
  \item
    What is the probability that the elderly employment rate is at least
    0.5 (50\%)?
  \item
    What is the probability that the elderly employment rate is between
    0.4 (40\%) and 0.6 (60\%)?
  \end{enumerate}
\end{enumerate}

\begin{enumerate}
\def\labelenumi{\arabic{enumi}.}
\setcounter{enumi}{4}
\item
  In any year, the weather can inflict storm damage to a home. From year
  to year, the damage is random. Let \(Y\) denote the dollar value of
  damage in any given year. Suppose that in 95\% of the years, there is
  no damage (\(Y=0\)), but that in 5\% of the years, \(Y = 20000\).

  \begin{enumerate}
  \def\labelenumii{\alph{enumii}.}
  \item
    What are the mean and standard deviation of the damage in any year?
  \item
    Consider an ``insurance pool'' of 100 people whose homes are
    sufficiently dispersed so that, in any year, the damange to
    different homes can be viewed as inddependently distributed random
    variables. Let \(\bar{Y}\) denote the average damage to these 100
    homes in a year (i) What is \(E[\bar{Y}]\)? (i) What is the
    probability that \(\bar{Y}\) exceeds \$2000?
  \end{enumerate}
\end{enumerate}

\begin{enumerate}
\def\labelenumi{\arabic{enumi}.}
\setcounter{enumi}{5}
\item
  Grades on a standardized test are known to have a mean of 1000 for
  students in the United States. The test is administered to 453
  randomly selected students in Florida; in this sample, the mean is
  1013 and the standard deviation (\(s\)) is 108.

  \begin{enumerate}
  \def\labelenumii{\alph{enumii}.}
  \item
    Construct a 95\% confidence interval for the average test score for
    Florida students
  \item
    Is there statistically significant evidence that Florida students
    perform differently than other students in the United States? How do
    you know?
  \end{enumerate}
\end{enumerate}

\textbf{For the following question, make sure you submit your log-file
alongside your answers!}

\begin{enumerate}
\def\labelenumi{\arabic{enumi}.}
\setcounter{enumi}{6}
\item
  Download \href{../countymurders.dta}{\texttt{countymurders.dta}} to
  answer this question.\footnote{Remember to move this file into your
    \href{../01-lab/\#working-directories-important}{working directory}!}
  The variable \(murders\) is the number of murders reported in the
  county. The variable \(execs\) is the number of executions that took
  place of people sentenced to death in the given country. Most states
  in the United States have the death penalty, but several do not.

  \begin{enumerate}
  \def\labelenumii{\alph{enumii}.}
  \item
    Keep only data from the year 1996. How many counties are there in
    the data set? Of these, how many have zero murders. What percentage
    of countries have zero executions?
  \item
    What is the largest number of murders in a county? What is the
    largest number of executions in a county?
  \item
    Compute the correlation coefficient \(r\) between \texttt{murders}
    and \texttt{execs} and describe what you find.\footnote{Remember,
      you can use \texttt{correlate\ var1\ var2} to look at the
      correlation between two variables.} Estimate the correlation
    coefficient between \texttt{murdrate} and \texttt{execrate}. Why do
    the two coefficients differ so much?
  \item
    What are two characteristics in the data that are highly correlated
    with county murder rates?\footnote{If you want to look at the
      correlation between lots of variables, you can use
      \texttt{correlate\ var1\ var2\ ...\ var99}. If you want to refer
      to a lot of variables, an asterisk (*) can act as a ``wild.'' So
      if you use \texttt{correlate\ var*}, you'll receive a correlation
      matrix of every variable with a name that starts with ``var.'' If
      you use \texttt{correlate\ *var*}, it will give you a correlation
      matrix of every variable with the letters ``var'' somewhere in the
      name.} What are their correlation coefficients?
  \item
    What is median real per-capita income?\footnote{\texttt{tabstat} is
      your friend!}
  \item
    Generate a variable, \texttt{highinc} that is equal to 1 if a county
    has above-median real per capita income, and 0 otherwise. What is
    \(E[rpcpersinc  | highinc =0]\)? What is
    \(E[rpcpersinc | highinc =1]\)?
  \item
    Consider a two-sided hypothesis test of whether murder rates are
    different between counties with high (above median) vs low (below
    median) real per-capita personal income. Assume the two samples are
    independent, with equal variances.

    \begin{enumerate}
    \def\labelenumiii{\alph{enumiii}.}
    \tightlist
    \item
      First, write a null and alternative hypothesis
    \item
      Use Stata to conduct the hypothesis test. What is the relevant
      t-statistic?\footnote{The help file for \texttt{ttest} will be
        useful. Here we are conducting a two-sample t-test using groups.
        You will want to use the \texttt{highinc} variable you generated
        earlier.}
    \item
      Can you reject the null hypothesis at the 5\% level?
    \end{enumerate}
  \item
    Generate a variable, \texttt{perc1029}, that is equal to the share
    of the population between the ages of 10 and 29. What is the median
    share of the population by county that is ages \textbf{10-29}?
  \item
    Generate a variable, \texttt{young}, that is equal to 1 if a county
    has an above-median share of the population that is age 10-29, and 0
    otherwise. What is \(E[perc1029| young = 0]\)? What is
    \(E[perc1029| young =1]\)?
  \item
    Consider a two-sided hypothesis test of whether murder rates are
    different between states with a high (above-median) share of the
    population ages \textbf{10-29} versus a low share. Assume the two
    samples are independent, with equal variances.

    \begin{enumerate}
    \def\labelenumiii{\alph{enumiii}.}
    \tightlist
    \item
      First, write a null and alternative hypothesis
    \item
      Use Stata to conduct the hypothesis test. What is the relevant
      t-statistic?
    \item
      Can you reject the null hypothesis at the 5\% level?
    \end{enumerate}
  \end{enumerate}
\end{enumerate}

\subsection*{Sources}\label{sources}
\addcontentsline{toc}{subsection}{Sources}

\href{../countymurders.dta}{\texttt{countymurders.dta}}

\emph{Source: Compiled by J. Monroe Gamble for a Summer Research
Opportunities Program (SROP) at Michigan State University, Summer 2014.
Monroe obtained data from the U.S. Census Bureau, the FBI Uniform Crime
Reports, and the Death Penalty Information Center.}

\end{document}
